\documentclass[12pt,a4paper,UTF8]{ctexart}
\usepackage[a4paper,top=22mm,bottom=22mm,left=22mm,right=22mm,headheight=15pt]{geometry}
\usepackage{amsmath,amssymb,mathtools,bm}
\usepackage{graphicx}
\usepackage{xcolor}
\usepackage{booktabs,longtable,array}
\usepackage{enumitem}
\usepackage{fancyhdr}
\usepackage{caption}
\usepackage{float}
\usepackage{microtype}
\usepackage{setspace}
\usepackage[hidelinks,unicode]{hyperref}
\usepackage[most]{tcolorbox}

\definecolor{ContestBlue}{HTML}{17365D}
\definecolor{SoftBlue}{HTML}{EEF4FA}
\definecolor{SoftGray}{HTML}{F5F6F7}
\hypersetup{pdftitle={第102套},pdfauthor={物理竞赛教学资料}}

\setstretch{1.24}
\setlength{\parindent}{2em}
\setlength{\parskip}{0.35em}
\setlist{itemsep=0.35em,topsep=0.4em}
\captionsetup{font=small,labelfont=bf,justification=centering}
\setcounter{secnumdepth}{0}

\pagestyle{fancy}
\fancyhf{}
\fancyhead[L]{\small 第102套}
\fancyfoot[C]{\thepage}
\renewcommand{\headrulewidth}{0.4pt}

\ctexset{
  section={format=\Large\bfseries\color{ContestBlue},beforeskip=1.2em,afterskip=0.7em},
  subsection={format=\large\bfseries\color{ContestBlue},beforeskip=1.0em,afterskip=0.5em},
  subsubsection={format=\normalsize\bfseries,beforeskip=0.8em,afterskip=0.35em}
}

\makeatletter
\newcommand{\tightlist}{\setlength{\itemsep}{0pt}\setlength{\parskip}{0pt}}
\newcommand{\maxwidth}{\ifdim\Gin@nat@width>\linewidth\linewidth\else\Gin@nat@width\fi}
\newcommand{\maxheight}{\ifdim\Gin@nat@height>0.72\textheight 0.72\textheight\else\Gin@nat@height\fi}
\newcommand{\pandocbounded}[1]{\begingroup\setkeys{Gin}{width=\maxwidth,height=\maxheight,keepaspectratio}#1\endgroup}
\makeatother

\begin{document}

\begin{titlepage}
  \centering
  \vspace*{58mm}
  {\Huge\bfseries\color{ContestBlue} 第102套\par}
  \vfill
  \vspace*{28mm}
\end{titlepage}


\section{第 1 题}\label{ux7b2c-1-ux9898}

白球为均匀球体，质量为 \(m\)。击球后，白球以初速度 \(v_0\)
撞向一个初始静止、质量也为
\(m\)、半径相同的黑球。由于瞄准位置不固定，两球碰撞时的碰撞参数可以任意变化。不计球与球、球与桌面之间的一切摩擦。

\begin{enumerate}
\def\labelenumi{\arabic{enumi}.}
\tightlist
\item
  若碰撞不是完全弹性的，恢复系数为
  \(e=0.5\)，求碰撞后白球可能达到的最大散射角
  \(\theta_{1\max}\)。散射角定义为白球末速度与初速度的夹角。
\item
  将白球换成半径相同、质量为 \(5m\) 的新白球。新白球仍以初速度 \(v_0\)
  撞向初始静止、质量为 \(m\)
  的黑球，且碰撞完全弹性。求碰撞后新白球可能达到的最大散射角
  \(\theta_{2\max}\)。
\end{enumerate}

\newpage

\section{第 2 题}\label{ux7b2c-2-ux9898}

一个半径为 \(R\)、质量为 \(M\) 的半圆柱放在水平面上。另一个半径为
\(r\)、质量为 \(m\)
的均匀实心小圆柱从半圆柱顶端由静止开始滚下。小圆柱与半圆柱之间的摩擦因数为
\(\mu_1\)，半圆柱与地面之间的摩擦因数为
\(\mu_2\)。小圆柱圆心相对半圆柱圆心转过的角度记为 \(\theta\)。

\begin{enumerate}
\def\labelenumi{\arabic{enumi}.}
\tightlist
\item
  \(\mu_1\)、\(\mu_2\) 均足够大。求小圆柱离开半圆柱时的角度 \(\theta\)。
\item
  \(\mu_2\) 足够大而 \(\mu_1\)
  较小。求小圆柱与半圆柱之间开始打滑时的角度 \(\theta\)。
\item
  \(\mu_1\) 足够大而 \(\mu_2\) 较小。求半圆柱与地面之间开始打滑的角度
  \(\theta\) 所满足的一元方程。
\item
  取 \(M=2m\)、\(\mu_2=0\)。当 \(\mu_1\)
  足够大时，判断小圆柱能否离开半圆柱；若能，求离开角度。当 \(\mu_1\)
  较小时，求小圆柱与半圆柱之间开始打滑的角度 \(\theta\)
  所满足的一元方程。
\end{enumerate}

\begin{figure}[H]
\centering
\includegraphics[width=0.9\linewidth,height=\textheight,keepaspectratio,alt={图 2-1　小圆柱沿半圆柱外表面滚下，图示角度从竖直方向量起。}]{../figures/question-02.pdf}
\caption*{图 2-1　小圆柱沿半圆柱外表面滚下，图示角度从竖直方向量起。}
\end{figure}

\newpage

\section{第 3 题}\label{ux7b2c-3-ux9898}

在 \(xOy\) 平面内有一半径为 \(r\)、圆心位于原点 \(O\)
的细圆环。圆环上一点的位置由其半径与 \(x\) 轴正方向的夹角 \(\theta\)
表示，其中 \(0\leq\theta\leq2\pi\)。

\begin{enumerate}
\def\labelenumi{\arabic{enumi}.}
\tightlist
\item
  圆环均匀带电，线电荷密度为常量 \(\lambda\)。

  \begin{enumerate}
  \def\labelenumii{\arabic{enumii}.}
  \tightlist
  \item
    求原点处的电势 \(\varphi_0\)。
  \item
    求原点附近三维空间中的电场，并保留到坐标小量的二阶精度。
  \end{enumerate}
\item
  圆环的线电荷密度为
  \[\lambda(\theta)=\lambda_0\frac{\theta}{2\pi},\qquad 0\leq\theta\leq2\pi.\]

  \begin{enumerate}
  \def\labelenumii{\arabic{enumii}.}
  \tightlist
  \item
    求原点处的电势 \(\varphi_1\) 和电场 \(\boldsymbol E_0\)。
  \item
    求原点附近三维空间中的电场，并保留到坐标小量的一阶。
  \end{enumerate}
\end{enumerate}

\newpage

\section{第 4 题}\label{ux7b2c-4-ux9898}

一个半径为 \(R\)、均匀带电量 \(+Q_0\)（\(Q_0>0\)）的绝缘球壳，以角速度

\[\boldsymbol\omega=-\omega\,\hat{\boldsymbol x}\]

匀速转动。

\begin{enumerate}
\def\labelenumi{\arabic{enumi}.}
\tightlist
\item
  电流回路的磁矩定义为
  \(\mathrm d\boldsymbol m=I\,\mathrm d\boldsymbol S\)，其中 \(I\)
  为电流，\(\mathrm d\boldsymbol S\)
  的大小为回路面积、方向由右手定则确定。求转动球壳的总磁矩
  \(\boldsymbol m_0\)。
\item
  写出电偶极子 \(\boldsymbol p\) 在球坐标 \((r,\theta)\)
  处所产生电场的径向分量 \(E_r\) 与极角分量 \(E_\theta\)。
\item
  利用对应关系
  \(1/(4\pi\varepsilon_0)\leftrightarrow\mu_0/(4\pi)\)、\(\boldsymbol p\leftrightarrow\boldsymbol m\)，求远区
  \(r\gg R\) 处的磁感应强度 \(\boldsymbol B\)。角度 \(\theta\) 从 \(+x\)
  轴量起。
\item
  在距球心 \(r\gg R\) 处放置电量为 \(-q\)（\(q>0\)）、质量为 \(m\)
  的点电荷。若它在 \(x=\text{常量}\) 的平面内绕 \(x\)
  轴做匀速圆周运动，求轨道极角 \(\theta\) 及角速度
  \(\omega'\)。忽略辐射、阻力和重力。
\end{enumerate}

\begin{figure}[H]
\centering
\includegraphics[width=0.9\linewidth,height=\textheight,keepaspectratio,alt={图 4-1　带电球壳绕负 x 方向转动，负点电荷在远区绕 x 轴做圆周运动。}]{../figures/question-04.pdf}
\caption*{图 4-1　带电球壳绕负 x 方向转动，负点电荷在远区绕 x
轴做圆周运动。}
\end{figure}

\newpage

\section{第 5 题}\label{ux7b2c-5-ux9898}

如图（a），光滑水平面上有两个均匀绝缘圆环，半径分别为 \(2R\) 和
\(R\)，质量均为 \(m\)。轻杆 \(O_1O_2\)
在两圆环的圆心处分别与圆环铰接，两圆环接触处始终无相对滑动。在小圆环
\(O_2\) 的右端固定电量为 \(q\) 的轻质点电荷。

\begin{enumerate}
\def\labelenumi{\arabic{enumi}.}
\tightlist
\item
  在以 \(O_1\) 为圆心、半径为 \(2R\) 的区域内加竖直方向的均匀磁场，且
  \(B=B_0+kt\)。固定圆环 \(O_1\)，将系统从图（a）位置静止释放。求杆
  \(O_1O_2\) 转过角度 \(\theta\) 时的角速度 \(\omega_1\) 和角加速度
  \(\beta_1\)。
\item
  接第（1）问，但允许圆环 \(O_1\) 绕过自身圆心的竖直轴自由转动。已知杆
  \(O_1O_2\) 转过角度 \(\theta\) 的过程用时为 \(t\)，求杆的角速度
  \(\omega_2\) 和角加速度 \(\beta_2\)。允许用点电荷相对 \(O_1\) 的方位角
  \(\theta'\) 及 \(\dot\theta'\) 表示结果，但必须给出 \(\theta'\)
  的表达式。
\item
  如图（b），固定圆环 \(O_1\)，将 \(O_2\) 换成半径为 \(2R\)、质量仍为
  \(m\) 的圆环 \(O_3\)。磁场区域扩展为以 \(O_1\) 为圆心、半径大于 \(6R\)
  的圆形区域，仍有 \(B=B_0+kt\)，两环接触处无相对滑动，点电荷仍为
  \(q\)。求系统从图（b）位置静止释放瞬间杆 \(O_1O_3\) 的角加速度
  \(\beta_3\)。
\item
  接第（3）问，求杆 \(O_1O_3\) 转过 \(180^\circ\) 时的角速度
  \(\omega_4\) 和角加速度 \(\beta_4\)。
\end{enumerate}

\begin{figure}[H]
\centering
\includegraphics[width=0.9\linewidth,height=\textheight,keepaspectratio,alt={图 5-1　两种外切圆环系统：（a）半径分别为 2R、R；（b）两环半径均为 2R。}]{../figures/question-05.pdf}
\caption*{图 5-1　两种外切圆环系统：（a）半径分别为
2R、R；（b）两环半径均为 2R。}
\end{figure}

\newpage

\section{第 6 题}\label{ux7b2c-6-ux9898}

一容器中装有质量 \(m_w=9\,\mathrm g\) 的水和物质的量
\(n_a=0.5\,\mathrm{mol}\) 的空气。初始时容器体积为
\(V_0=61.4\,\mathrm L\)。随后绝热压缩到体积
\(V_1\)，并认为该过程中水全部处于气态。压缩后气体温度为
\(T_1=373\,\mathrm K\)、总压强为 \(p_1=1\,\mathrm{atm}\)，其中
\(1\,\mathrm{atm}=1.01\times10^5\,\mathrm{Pa}\)。再使系统在 \(T_1\)
下等温压缩到体积 \(V_1/4\)，并认为系统中有足够的凝结核。

本题条件下，水的汽化热为
\(L=2.26\times10^5\,\mathrm{J\,kg^{-1}}\)，液态水体积可忽略。气体常量为

\[R=8.31\,\mathrm{J\,mol^{-1}\,K^{-1}}.\]

水蒸气和空气分别视为多原子分子、双原子分子理想气体，只有平动和转动自由度被激发。

\begin{enumerate}
\def\labelenumi{\arabic{enumi}.}
\tightlist
\item
  求绝热压缩前水蒸气与空气的分压 \(p_{w0}\)、\(p_{a0}\)。
\item
  求随后等温压缩过程中系统放出的热量。
\end{enumerate}

\newpage

\section{第 7 题}\label{ux7b2c-7-ux9898}

在惯性系 \(S\) 中有一半径为 \(R\) 的圆 \(C\)，圆周上有一点 \(O\)。在
\(t=0\) 时，质点 \(A\)、\(B\) 同时从 \(O\) 出发，均以恒定速率
\(\beta c\) 运动。\(A\) 沿圆周运动，\(B\) 沿圆在 \(O\)
点的切线做匀速直线运动，初始时二者速度同向。

\(A\)、\(B\) 各携带一只钟。\(B\) 的钟指示 \(B\) 参考系中的时间，并在
\(O\) 点调为零；\(A\) 的钟与静止在 \(S\) 系、指示 \(S\) 系时间的钟在
\(S\) 系中同步。

\begin{enumerate}
\def\labelenumi{\arabic{enumi}.}
\tightlist
\item
  求当 \(B\) 认为自己的钟指示 \(t_B\) 时，\(B\) 认为 \(A\) 的钟示数
  \(t_A\) 所满足的方程，不必求解。
\item
  求当 \(A\) 认为自己的钟指示 \(t_A\) 时，\(A\) 在其瞬时惯性系中认为
  \(B\) 的钟示数 \(t_B\)。
\end{enumerate}

\begin{figure}[H]
\centering
\includegraphics[width=0.844\linewidth,height=\textheight,keepaspectratio,alt={图 7-1　A 沿圆周运动，B 沿 O 点切线运动，二者初速度同向。}]{../figures/question-07.pdf}
\caption*{图 7-1　A 沿圆周运动，B 沿 O 点切线运动，二者初速度同向。}
\end{figure}

\newpage

\section{第 8 题}\label{ux7b2c-8-ux9898}

一个能量为 \(E_\gamma=100.0\,\mathrm{MeV}\)
的光子撞击一个初始静止、质量为 \(m_n=3.932\,\mathrm{GeV}/c^2\)
的原子核。光子被吸收后产生一对正、负电子，原子核发生反冲，反冲动量大小记为
\(p_R\)。电子静能为 \(m_ec^2=0.51\,\mathrm{MeV}\)。

\begin{enumerate}
\def\labelenumi{\arabic{enumi}.}
\tightlist
\item
  在满足能量---动量守恒的前提下，求反冲核动量 \(p_R\) 的最小值。
\item
  求反冲核速度方向与入射光子方向的最大夹角 \(\theta_M\)。
\end{enumerate}

\end{document}
